\section{Les évènements dans les règles Jerboa }
\label{sec:formalisation-evenement-regles-jerboa}

\highlight{chapeau}

%\maxime { adapter les propositions aux schémas: wip}

\subsection{Création}

\begin{figure}[htbp!]
	\centering
	\includegraphics[height=50pt]{rule-face-triangulation-3D-no-color}
	\caption{Règle de triangulation de face}
	\label{fig:formalisation-evenement-regles-jerboa-rule-face-triangulation}
\end{figure}

Continuons avec la règle de triangulation, cette fois sous la forme d'une règle
Jerboa
(Figure~\ref{fig:formalisation-evenement-regles-jerboa-rule-face-triangulation}). %
Dans la section~\ref{sec:formalisation-evenement-regles-transfo-graphe-creation}
(page~\pageref{sec:formalisation-evenement-regles-transfo-graphe-creation}), nous avons étudié la création d'un sommet dans une face carrée suite à sa transformation directe par la règle instanciée de triangulation (Figure~\ref{fig:formalisation-evenement-regles-transfo-graphe-triangulation} page~\pageref{fig:formalisation-evenement-regles-jerboa-rule-face-triangulation}). %
Nous allons maintenant montrer que la création du sommet $R\orbit{1,2,3}(n2)$
dans la règle Jerboa détermine, elle aussi, la création d'un sommet dans la
transformation directe. %

\begin{proposition}[Création d'orbite]\label{prop:formalisation-evenement-regles-jerboa-creation}
	Soit une règle Jerboa $r~: L \longrightarrow R$, un nœud d'accroche $h$ de $L$, un morphisme de
	filtrage $\underline{m}~: \underline{L} \rightarrow G$ sur une G-carte $G$, un type
	d'orbite $\orbit{\omega}$ et $v$ un nœud de $R$.

	L'orbite $\underline{R}\orbit{\omega}(v)$ est créée dans la
	règle si et seulement si toutes les orbites $H\orbit{\omega}(m((v,v')))$ sont créées
	dans la transformation directe $G \xlongrightarrow{r, \underline{m}} H$, pour tous les brins
	$v'$ de l'orbite d'instanciation $O = G\orbit{l(h)}(\underline{m}(h))$.
\end{proposition}

\begin{proof}
	Soit $v$ un nœud de $R$ tel que l'orbite $\underline{R}\orbit{\omega}(v)$ est créée
	par la règle, c'est-à-dire tel que
	$\underline{R}\orbit{\omega}(v) \cap \underline{L} = \emptyset$. Soit
	$O = G\orbit{l(h)}(\underline{m}(h))$, l'orbite d'instanciation.

	Soit $w$ un nœud quelconque de l'orbite $\underline{R}\orbit{\omega}(v)$. $w$ est un
	nœud créé de la règle (c'est-à-dire de $R \setminus L$). %
	Donc, par construction de la règle instanciée,
	tous les nœuds $(w,v')$ pour tous les brins $v'$ de l'orbite d'instanciation
	$O$, sont des nœuds créés dans la règle instanciée $r(\underline{m})$.

	Donc toutes les orbites $R(O)\orbit{\omega}((v,v'))$ pour tous les brins $v'$ de
	l'orbite d'instanciation $O$ sont des orbites créées dans la règle instanciée
	$r(\underline{m})$. Et donc, d'après la proposition \ref{prop:base-rule-create},
	toutes les orbites $H\orbit{\omega}(m((v,v')))$ sont créés dans la transformation
	directe $G \xlongrightarrow{r, \underline{m}} H$, pour tous les brins $v'$ de l'orbite
	d'instanciation $O$.
	\\

	Réciproquement, supposons que toutes les orbites $H\orbit{\omega}(m'((v,v')))$ sont créés dans la transformation directe $G \xlongrightarrow{r, \underline{m}} H$, pour tous les brins $v'$ de l'orbite d'instanciation $O$.
	Donc, d'après la proposition \ref{prop:base-rule-create}, toutes les orbites $R(O)\orbit{\omega}((v,v'))$ pour tous les brins $v'$ de l'orbite d'instanciation $O$ sont des orbites créées dans la règle instanciée $r(\underline{m})$.
	Donc par construction de l'instanciation, $\underline{R}\orbit{\omega}(v)$ est une orbite créée dans la règle $r$.
\end{proof}

\begin{figure}[ht]
	\begin{subfigure}{.5\linewidth}
		\centering
		\includegraphics[height=100pt]{g-carte-triangulation-instanciation-carre-G}
		\caption{Instanciation d'une face de \(G_{quad}\)}
		\label{fig:formalisation-evenement-regles-jerboa-instanciation-quad-triangulation-G}
	\end{subfigure}
	\begin{subfigure}{.5\linewidth}
		\centering
		\includegraphics[height=100pt]{g-carte-triangulation-instanciation-carre-H}
		\caption{Instanciation d'une face de \(H_{quad}\)}
		\label{fig:formalisation-evenement-regles-jerboa-instanciation-quad-triangulation-H}
	\end{subfigure}
	\begin{subfigure}{.5\linewidth}
		\centering
		\includegraphics[height=100pt]{g-carte-triangulation-instanciation-triangle-G}
		\caption{Instanciation d'une face de \(G_{tri}\)}
		\label{fig:formalisation-evenement-regles-jerboa-instanciation-tri-triangulation-G}
	\end{subfigure}
	\begin{subfigure}{.5\linewidth}
		\centering
		\includegraphics[height=100pt]{g-carte-triangulation-instanciation-triangle-H}
		\caption{Instanciation d'une face de \(H_{tri}\)}
		\label{fig:formalisation-evenement-regles-jerboa-instanciation-tri-triangulation-H}
	\end{subfigure}
	\caption{Instanciations de la règle de triangulation pour différents polygones}
	\label{fig:formalisation-evenement-regles-jerboa-inst-quad-triangulation}
\end{figure}

\begin{figure}[htb!]
	\begin{subfigure}{.5\textwidth}
		\centering
		\includegraphics[width=120pt]{g-carte-triangulation-face-triangle-G}
		\caption{\(G_{tri}\)}
		\label{fig:formalisation-evenements-regle-jerboa-creation-triangulation-tri-G}
	\end{subfigure}
	\begin{subfigure}{.5\textwidth}
		\centering
		\includegraphics[width=120pt]{g-carte-triangulation-face-triangle-H}
		\caption{\(H_{tri}\)}
		\label{fig:formalisation-evenements-regle-jerboa-creation-triangulation-tri-H}
	\end{subfigure}
	\caption{Triangulation d'une face triangulaire}
	\label{fig:formalisation-evenements-regle-jerboa-creation-triangulation-tri}
\end{figure}
% Création de l'orbite dans la règle
Dans la règle de triangulation
(Figure~\ref{fig:formalisation-evenement-regles-jerboa-rule-face-triangulation}),
l'orbite \(R\orbit{1,2,3}(n2)\) contient uniquement le n\oe ud \(n2\). %
Puisque celui-ci n'appartient pas à \(L\), nous pouvons en déduire qu'il est
créé et donc l'orbite est créée dans \(R\). %
% Création de l'orbite dans la règle instanciée, on reprendre la quad
Dans la transformation directe, pour une face carrée, la règle est instanciée
sous la forme de la
règle~\ref{fig:formalisation-evenement-regles-transfo-graphe-triangulation}. %
Nous avons montré dans la section précédente que l'application de cette règle
instanciée sur une face carrée de \(G\) crée bien un sommet dans \(H\). %
% Création du sommet dans les faces instanciées
Toutefois, ici, cette règle peut être appliquée sur tout type de face peu
importe le nombre de côtés comme une face carrée
(Figures~\ref{fig:formalisation-evenement-regles-jerboa-instanciation-quad-triangulation-G}\ref{fig:formalisation-evenement-regles-jerboa-instanciation-quad-triangulation-H})
ou
triangulaire~\ref{fig:formalisation-evenement-regles-jerboa-instanciation-tri-triangulation-G}\ref{fig:formalisation-evenement-regles-jerboa-instanciation-tri-triangulation-H}. %
Pour chacune des deux instanciations, le sommet est bien créé
\(M(G_{quad}\orbit{1,2,3}(a_{2}))\) et \(M(G_{tri}\orbit{1,2,3}(m_{2}))\) ainsi
que dans la transformation directe avec, par exemple, la face
triangulaire (Figures~\ref{fig:formalisation-evenement-regles-jerboa-instanciation-tri-triangulation-H}
et \ref{fig:formalisation-evenement-regles-jerboa-inst-quad-triangulation}). %
% Création du sommet dans Hquad et Htri
Enfin, puisque le sommet est créé dans la règle Jerboa et ses instanciations,
celui-ci est créé dans \(H_{tri}\)
(Figures~\ref{fig:formalisation-evenements-regle-jerboa-creation-triangulation-tri-G}
et \ref{fig:formalisation-evenements-regle-jerboa-creation-triangulation-tri-H})

~

L'analyse des règles Jerboa nous permet de définir ce que nous appelons les
orbites d'origines. %
Pour le présenter simplement, une orbite d'origine est l'ancêtre invariant d'une
entité topologique. %

\begin{definition}[Orbite d'origine]
	Soit une règle $r~: L \longrightarrow R$, $n$ un nœud de $R$ et $\orbit{\omega}$ un type d'orbite.

	Une \emph{orbite d'origine} de $\underline{R}\orbit{\omega}(n)$ dans $L$ est une orbite $\underline{L}\orbit{\omega'}(n')$ où~:
	\begin{itemize}
		\item $n'$ est un nœud de $L$,

		\item $\orbit{\omega'}$ est l'ensemble des dimensions  $i$, où $i$ est renommé dans un arc implicite de $n'$ ; c'est-à-dire que $(l_L(n') \mapsto l_R(n))(i) = j \in \orbit{\omega}$

	\end{itemize}
	Soit un morphisme de filtrage $\underline{m} : \underline{L} \rightarrow G$ sur une G-carte $G$, $G \xlongrightarrow{r, \underline{m}} H$ la transformation directe correspondante et $O$ l'orbite d'instanciation. Soit $v'$ un nœud de $O$.

	Une {\em origine} de l'orbite $R(O)\orbit{\omega}((n,v'))$ est une orbite $L(O)\orbit{\omega'}((n',v'))$.
\end{definition}

Dans notre exemple, le sommet créé est le dual de la face triangulée et donc,
intuitivement, la face est l'origine du sommet. %
Plus formellement, dans l'orbite \(R\orbit{1,2,3}(n2)\), l'étiquette \(1\) (resp.)
\(2\) et \(3\)) est une réécriture de l'étiquette \(0\) (resp. \(1\) et \(3\))
du n\oe ud \(n0\) dans $L$. %
L'ensemble des étiquettes réécrites désignent alors une face \(\orbit{0,1,3}\)
comme origine du sommet. %

\subsection{Scission}

% Reprenons la règle de triangulation
% (Figure~\ref{fig:formalisation-evenement-regles-jerboa-rule-face-triangulation}) et son instanciation (Figure~\ref{fig:formalisation-evenement-regles-jerboa-inst-quad-triangulation}). %
% La règle filtre une face $L\orbit{0,1,3}(n0)$ puis la scinde en plusieurs faces triangulaires. %
% On peut observer, dans la règle, que l'orbite $L\orbit{0,1,3}(n0)$ présente
% toutes les étiquettes de $\orbit{0,1,3}$, on s'intéresse particulièrement aux
% $1$-étiquettes puisque la face filtrée est scindée au niveau de ses
% sommets. %
% Dans l'orbite $R\orbit{0,1,3}(n0)$, aucun des n\oe uds ne présente de
% $1$-étiquette. %
% Dans cette règle, cette réécriture d'étiquettes caractérise la scission d'une face.
Dans la
section~\ref{sec:formalisation-evenement-regles-transfo-graphe-scission} (page~\pageref{sec:formalisation-evenement-regles-transfo-graphe-scission}), nous avons étudié la scission d'une face suite à sa transformation directe par la règle instanciée de triangulation. %
Nous allons maintenant montrer que la scission de la face \(L\orbit{0,1,3}(n0)\)
dans la règle Jerboa détermine, elle aussi, la scission d'une face dans la
transformation directe. %

\begin{proposition}[Scission d'orbite]\label{prop:formalisation-evenement-regles-jerboa-scission}
	Soit une règle $r~: L \longrightarrow R$, un nœud d'accroche $h$ de $L$, un morphisme de
	filtrage $\underline{m}~: \underline{L} \rightarrow R$ sur une G-carte $G$, un type d'orbite
	$\orbit{\omega}$, deux nœuds préservés $v_1$ et $v_2$ de $\underline{K} = \underline{L} \cap \underline{R}$. Notons $O = G\orbit{l(h)}(\underline{m}(h))$ l'orbite d'instanciation.

	$\underline{L}\orbit{\omega}(v_1)$ est scindé dans la règle si :
	\begin{itemize}
		\item la scission est due à des arcs explicites, c'est-à-dire si $\underline{L}\orbit{\omega}(v_1)$ est scindée dans $\underline{R}$~;

		\item la scission est due à des arcs implicites, c'est-à-dire si il existe un nœud $v$ de $\underline{L}\orbit{\omega}(v_1)$ dont l'étiquette $l(v)$ a pour k-ème position une dimension de $\orbit{\omega}$, et que pour tous les nœuds $w$ de $\underline{R}\orbit{\omega}(v_1)$ l'étiquette de $l(w)$ a pour k-ème position est un "$\_$" ou une dimension qui n'appartient pas à $\orbit{\omega}$.
	\end{itemize}

	\begin{enumerate}

		\item
		      Si $\underline{L}\orbit{\omega}(v_1)$ est scindée en deux orbites $\underline{R}\orbit{\omega}(v_1)$ et $\underline{R}\orbit{\omega}(v_2)$ dans la règle du à des arcs explicites.
		      Et si l'orbite $\underline{R}\orbit{\omega}(v_1)$ est complète. %
		      Alors, pour tout brin $v'$ de l'orbite d'instanciation $O$, l'orbite $G\orbit{\omega}(m((v_1,v')))$ est scindée en deux orbites $H\orbit{\omega}(m'((v_1, v')))$ et $H\orbit{\omega}(m'((v_2, v')))$ dans la transformation directe $G \xlongrightarrow{r,\underline{m}} H$ .

		\item Si l'orbite $\underline{L}\orbit{\omega}(v_1)$ est scindée en deux orbites $\underline{R}\orbit{\omega}(v_1)$ et
		      $\underline{R}\orbit{\omega}(v_2)$ dans la règle due à un arc explicite.
		      Si les orbites $\underline{R}\orbit{\omega}(v_1)$ et $\underline{R}\orbit{\omega}(v_2)$
		      sont toutes deux incomplètes.
		      Alors, pour chaque brin $v'$ de l'orbite d'instanciation $O$, soit $G\orbit{\omega}(m((v_1,v')))$ est scindée en deux orbites
		      $H\orbit{\omega}(m'((v_1, v')))$ et $H\orbit{\omega}(m'((v_2,v')))$, soit l'orbite est modifiée, dans la transformation directe
		      $G \xlongrightarrow{r, \underline{m}} H$.

		\item Si l'orbite $\underline{L}\orbit{\omega}(v_1)$ est scindée dans la règle du au $k$-ième arc implicite.
		      Alors pour tout brin $v'_1$ de l'orbite d'instanciation $O$ et $v'_2$ le $i$-voisin de $v'_1$ dans $O$ avec $i$ la dimension à la $k$-ième position de $l(h)$, $G\orbit{\omega}(m((v_1,v'_1)))$ est soit scindée en deux orbites
		      $H\orbit{\omega}(m'((v_1, v'_1)))$ et $H\orbit{\omega}(m'((v_1,v'_2)))$, soit modifiée en une orbite unique $H\orbit{\omega}(m'((v_1, v'_1))) = H\orbit{\omega}(m((v_1,v'_2)))$ dans la transformation directe
		      $G \xlongrightarrow{r, \underline{m}} H$.

		\item Réciproquement,

		      si l'orbite $G\orbit{\omega}(m((v_1, v'_1))$ est scindée en deux orbites
		      $H\orbit{\omega}(m'((v_1, v'_1)))$ et $H\orbit{\omega}(m'((v_1,v'_2)))$ alors il existe deux n\oe uds préservés $v_3$ et $v_4$ (de $\underline{K} = \underline{L} \cap \underline{R}$) tel que $m((v_3, v'_1)) \in G\orbit{\omega}(m((v_1, v'_1))$ et $m((v_4, v'_1)) \in G\orbit{\omega}(m((v_1, v'_1))$ et $\underline{L}\orbit{\omega}(v_3)$ est scindée en $\underline{R}\orbit{\omega}(v_3)$ et $\underline{R}\orbit{\omega}(v_4)$. %

	\end{enumerate}

\end{proposition}

\begin{proof}
	Soit deux nœuds préservés $v_1$ et $v_2$ de $\underline{K} = \underline{L} \cap \underline{R}$
	\begin{enumerate}
		\item Supposons que $\underline{L}\orbit{\omega}(v_1)$ est scindée en deux
		      orbites $\underline{R}\orbit{\omega}(v_1)$ et $\underline{R}\orbit{\omega}(v_2)$, avec au moins l'une des deux complète, dans la
		      règle du à des arcs explicites. Alors, par définition, il existe un chemin
		      $e_1 e_2 ... e_j$ dans $\underline{L}$ de source $v_1$, de cible $v_2$ et étiqueté par
		      $i_1 i_2 ... i_j$ sur $\orbit{\omega}$, mais il n'existe pas de tel chemin
		      $f_1 f_2 ... f_k$ dans $R$ de source $v_1$, de cible $v_2$ et étiqueté sur
		      $\orbit{\omega}$.

		      Donc dans la règle instanciée, pour tout nœud $v'$ de l'orbite
		      d'instanciation $O = G\orbit{l(h)}(\underline{m}(h))$, il existe un chemin
		      $(e_1,v') (e_2,v') ... (e_j,v')$ dans le motif gauche de la règle instanciée
		      $L(O)$ dont la source est $(s(e_1),v') = (v_1, v')$, la cible est
		      $(t(e_j), v') = (v_2, v')$ et étiqueté par le même mot $i_1 i_2 ... i_j$ sur
		      $\orbit{\omega}$.

		      Raisonnons par l'absurde. %
		      Supposons que pour un nœud $v'$ de l'orbite d'instanciation $O$, il
		      existe un chemin $(f_1,v') (f_2,v') ... (f_k,v')$ dans le motif droit de la
		      règle instanciée $R(O)$ dont la source est $(s(f_1),v') = (v_1, v')$, la cible
		      est $(t(e_j), v') = (v_2, v')$ et étiqueté sur $\orbit{\omega}$. Alors, par
		      construction de l'instanciation, le chemin $f_1 f_2 ... f_k$ dans $R$ de source
		      $v_1$, de cible $v_2$ et étiqueté sur $\orbit{\omega}$ existe dans la règle de
		      départ. Ce qui est contradictoire avec l'hypothèse. Donc il n'existe pas de
		      chemin $(f_1,v') (f_2,v') ... (f_k,v')$ dans le motif droit de la règle
		      instanciée ayant les propriétés attendues. Donc $R(O)\orbit{\omega}((v_1,v'))$ et
		      $R(O)\orbit{\omega}((v_2,v'))$ sont deux $\orbit{\omega}$-orbites distinctes.

		      Donc d'après la proposition \ref{prop:base-rule_split}, si
		      $\underline{R}\orbit{\omega}(v_1)$ et donc $R(O)\orbit{\omega}((v_1, v'))$ est entièrement
		      filtrée, alors l'orbite $G\orbit{\omega}(m((v_1,v')))$ est scindée en deux orbites
		      $H\orbit{\omega}(m'((v_1, v')))$ et $H\orbit{\omega}(m'((v_2, v')))$ dans la
		      transformation directe $G \xlongrightarrow{r,\underline{m}} H$ pour tout nœud
		      $v'$ de l'orbite d'instanciation $O$. %

		\item
		      Supposons que $\underline{L}\orbit{\omega}(v_1)$ est scindée en deux orbites incomplètes
		      $\underline{R}\orbit{\omega}(v_1)$ et $\underline{R}\orbit{\omega}(v_2)$ dans la règle du
		      à des arcs explicites.

		      Pour les mêmes raisons que dans le cas 1, $R(O)\orbit{\omega}((v_1,v'))$ et
		      $R(O)\orbit{\omega}((v_2,v'))$ sont deux $\orbit{\omega}$-orbites distinctes dans la
		      règles instanciée, pour tout nœud $v'$ de l'orbite d'instanciation $O$.

		      Donc d'après la proposition \ref{prop:base-rule_split}, si $\underline{R}\orbit{\omega}(v_1)$
		      et donc $R(O)\orbit{\omega}((v_1, v'))$ est incomplète, alors, pour tout nœud
		      $v'$ de l'orbite d'instanciation $O$, l'orbite
		      $G\orbit{\omega}(m((v_1,v')))$ est soit scindée en deux orbites
		      $H\orbit{\omega}(m'((v_1, v')))$ et $H\orbit{\omega}(m'((v_2, v')))$, soit modifiée dans
		      la transformation directe $G \xlongrightarrow{r,\underline{m}} H$.  %

		\item
		      Supposons que $\underline{L}\orbit{\omega}(v_1)$ est scindée due au $k$-ième arc
		      implicite. Soit $v'_1$ un nœud de l'orbite d'instanciation
		      $O = G\orbit{l(h)}(\underline{m}(h))$ et $v'_2$ le $i$-voisin de $v'_1$ dans
		      $O$, avec $i$ la $k$-ième dimension de l'étiquette $l(h)$ du nœud d'accroche $h$
		      de $L$.

		      Par hypothèse, l'un des nœuds $v$ de $\underline{L}\orbit{\omega}(v_1)$ a une
		      $k$-ième dimension $j$ dans l'$\orbit{\omega}$-orbite ($j \in \orbit{\omega}$), alors que tous
		      des n\oe uds de $\underline{R}\orbit{\omega}(v_1)$ ont leur $k$-ième dimension hors de
		      $\orbit{\omega}$. Donc les nœuds $(v, v'_1)$ et $(v, v'_2)$ sont $j$-voisins dans
		      l'instanciation $L(O)$, donc $(v, v'_1)$ et $(v, v'_2)$ sont dans la même
		      $\orbit{\omega}$-orbite dans $L(O)$, et donc $(v_1, v'_1)$ et $(v_2, v'_2)$ aussi.

		      Nous avons plusieurs cas :
		      \begin{itemize}
			      \item si $(v_1, v'_1)$ et $(v_2, v'_2)$ sont dans deux $\orbit{\omega}$-orbites
			            différentes de $R(O)$, d'après la proposition \ref{prop:base-rule_split}, selon
			            les cas, $H\orbit{\omega}(m'((v_1, v'_1)))$ et $H\orbit{\omega}(m'((v_1,v'_2)))$ sont
			            deux orbites différentes ou non de $H$, avec $H$ le résultat de la
			            transformation directe $G \xlongrightarrow{r, \underline{m}} H$ ; autrement dit
			            $G\orbit{\omega}(m((v_1, v'_1))$ est scindée en $H\orbit{\omega}(m'((v_1, v'_1)))$ et
			            $H\orbit{\omega}(m'((v_1,v'_2)))$ ou juste modifiée par la transformation directe
			            $G \xlongrightarrow{r, \underline{m}} H$ ;

			      \item si au contraire $(v_1, v'_1)$ et $(v_2, v'_2)$ sont dans la même
			            $\orbit{\omega}$-orbites de $R(O)$ (car il existe un chemin $e_1 e_2 ... e_l$ de $O$
			            de source $v'_1$ et de cible $v'_2$ qui ne comprend pas l'étiquette $j$, ce
			            chemin est réécrit en un chemin $(w, e_1) (w,e_2) ... (w, e_l)$ de $R(O)$ dont
			            toutes les étiquettes sont dans $\orbit{\omega}$), d'après la proposition
			            \ref{prop:base-rule-modify}, alors $G\orbit{\omega}(m((v_1, v'_1))$ est modifiée par
			            la transformation directe $G \xlongrightarrow{r, \underline{m}} H$. %
		      \end{itemize}
		\item
		      Soient $m((v_1,v'_1))$ et $m((v_2,v'_2))$ deux brins de $G$ filtrés par la règle avec $v_1$ et $v_2$ deux n\oe uds de $\underline{K} = \underline{L} \cap \underline{R}$ et $v'_1$ et $v'_2$ deux brins de $O$. %
		      Supposons que l'orbite $G\orbit{\omega}(m((v_1,v')))$ est scindée en deux orbites $H\orbit{\omega}(m'((v_1, v'_1)))$ et $H\orbit{\omega}(m'((v_2, v'_2)))$ dans la transformation directe. %
		      D'après la proposition~\ref{prop:base-rule_split}, il existe deux n\oe uds $k_1$ et $k_2$ préservés de la règle instanciée (de $K(O) = L(O) \cap R(O)$) tel que $m(k_1) \in G\orbit{\omega}(m((v_1, v'_1))$ et $m(k_2) \in G\orbit{\omega}(m((v_2,v'_2)$ et $L(O)\orbit{\omega}(k_1)$ est scindée en $R(O)\orbit{\omega}(k_1)$ et $R(O)\orbit{\omega}(k_2)$. %
		      Notons $k_1 = (v_3, v'_3)$ et $k_2 = (v_4, v'_4)$ avec $v_3$ et $v_4$ deux n\oe uds préservés de $\underline{K} = \underline{L} \cap \underline{R}$ et $v'_3$ et $v'_4$ deux brins de $O$. %
		      Nous avons donc $L(O)\orbit{\omega}((v_3,v'_3))$ scindée en $R(O)\orbit{\omega}((v_3, v'_3))$ et $R(O)\orbit{\omega}((v_4, v'_4))$. %
		      Donc il existe au moins chemin dans $L(O)$ entre $(v_3, v'_3)$ et $(v_4, v'_4)$, mais aucun chemin dans $R(O)$ entre $(v_3, v'_3)$ et $(v_4, v'_4)$. %
		      Donc, par construction de l'instanciation $L(O)$, il existe un chemin entre $(v_3, v'_3)$ et $(v_4, v'_3)$ (dont tous les n\oe uds sont de la forme $(v_i, v'_3)$) et un chemin entre $(v_4, v'_3)$ et $(v_4, v'_4)$ (dont tous les n\oe uds sont de la forme $(v_4, v'_i)$). %

		      Raisonnons par l'absurde. %
		      Supposons que la règle Jerboa ne comporte aucune scission. %
		      Par construction de l'instanciation~:
		      \begin{itemize}
			      \item Si la règle Jerboa n'a pas de scission explicite entre $v_3$ et $v_4$, c'est-à-dire que s'il existe un chemin entre $v_3$ et $v_4$ dans $L$ alors il en existe un dans $R$. %
			            Alors, d'après l'instanciation, s'il existe un chemin de $(v_3, v'_3)$ vers $(v_4, v'_3)$ dans $L(O)$, il en existe un dans $R(O)$. %

			      \item Si la règle Jerboa n'a pas de scission implicite entre $\underline{L}\orbit{\omega}(v_3)$ et $\underline{R}\orbit{\omega}(v_3)$. %
			            Alors, d'après l'instanciation, s'il existe un chemin de $(v_4, v'_3)$ vers $(v_4, v'_4)$ dans $L(O)$, il en existe un dans $R(O)$. %
		      \end{itemize}
		      Donc il existe un chemin dans $R(O)$ de $(v_3, v'_3)$ vers $(v_4, v'_4)$, ce qui est contradictoire avec l'hypothèse. %
		      Donc, la règle Jerboa comprend au moins une scission implicite ou explicite. %
	\end{enumerate}
\end{proof}

\begin{figure}[htbp!]
	\centering
	\includegraphics[height=50pt]{rule-unsew-a3}
	\caption{Règle de décousure de deux faces liées par des $3$-liaisons}
	\label{fig:formalisation-evenement-regles-jerboa-rule-unsew-a3}
\end{figure}

\begin{figure}[htbp!]
	\centering
	\begin{subfigure}{.45\textwidth}
		\centering
		\includegraphics[width=.85\textwidth]{g-carte-decousure-3-G}
		\caption{\label{fig:formalisation-evenements-regles-jerboa-evenement-scission-explicite-g}$G$}
	\end{subfigure}
	\begin{subfigure}{.45\textwidth}
		\centering
		\includegraphics[width=.85\textwidth]{g-carte-decousure-3-H}
		\caption{\label{fig:formalisation-evenement-regles-jerboa-evenement-scission-explicite-face-h}$H$}
	\end{subfigure}
	\caption{Décousure des $3$-liaisons dans une face}
	\label{fig:formalisation-evenement-regle-evenement-scission-explicite}
\end{figure}

%% 1/ Orbites complètes scission explicite
Pour illustrer le premier cas, une scission explicite où il existe au moins une
\(\orbit{\omega}\)-orbite complète, étudions la règle de décousure des \(3\)-arcs
d'une face. %
L'orbite \(L\orbit{0,1,3}(n0)\) est complète, de même, les deux orbites scindées
sont, elles aussi, complètes. %
Puisque dans la règle au moins l'une des \(\orbit{0,1,3}\)-orbites est complète,
alors tous les chemins entre les deux faces à séparer sont représentés. %
Donc la scission du \(3\)-arc explicite dans la règle provoque, comme nous
pouvons le voir dans la transformation
directe(Figure~\ref{fig:formalisation-evenement-regle-evenement-scission-explicite}),
la scission de toutes les \(3\)-liaisons entre les deux faces. %

%% 2/ Orbites incomplètes
\begin{figure}[htb!]
	\centering
	\includegraphics[height=50pt]{rule-insert-edge}
	\caption{Règle d'insertion d'arête entre deux sommets}
	\label{fig:formalisation-evenement-regles-jerboa-scission-rule-insert-edge}
\end{figure}

Pour le deuxième cas, où il n'existe pas d'\(\orbit{\omega}\)-orbite complète,
prenons la règle d'insertion d'arête
(Figure~\ref{fig:formalisation-evenement-regles-jerboa-scission-rule-insert-edge}). %
Une scission est détectés dans l'orbite \(R\orbit{0,1,3}(n0)\) due à la
suppression des \(1\)-arcs implicites. %
Toutefois, dans cette règle, il n'existe aucune \(\orbit{0,1,3}\)-orbite
complète. %
Donc la scission n'est pas certaines. %
Dans ce cas précis, il ne peut pas y avoir de scission si les orbites filtrées
appartiennent à deux faces distinctes. %
Le résultat de la transformation directe produit alors le résultat vu dans la
section~\ref{sec:formalisation-evenement-regles-transfo-graphe-scission} avec
l'exemple du cylindre
(Figures~\ref{fig:formalisation-evenement-regles-transfo-graphe-scission-cylindre-g}\ref{fig:formalisation-evenement-regles-transfo-graphe-scission-cylindre-h}). %



%% 3a/ Orbite complète scission implicite
Pour le troisième cas, commençons avec la règle de triangulation. %
La scission est due aux arcs implicites et plus particulièrement à la réécriture
des \(1\)-arcs. %
Dans chaque n\oe ud de \(R\orbit{0,1,3}(n0)\), les arcs implicites sont réécris à
partir du \(1\)-arc de \(L\orbit{0,1,3}(n0)\) sont soit supprimés, soit réécris
en \(2\)-arcs. %
Dans les instanciations de la règle
(Figure~\ref{fig:formalisation-evenement-regles-jerboa-inst-quad-triangulation}),
nous pouvons voir que les arêtes des faces filtrées sont séparées par la
suppression des \(1\)-liaisons et que les faces issues de la scission sont
elles-mêmes distinctes, car il n'existe pas de \(1\)-chemin qui connecte deux
faces. %

%% 3b/ orbite complète scission implicite => idempotence
\begin{figure}[htb!]
	\centering
	\includegraphics[height=50pt]{rule-unsew-volume-a2}
	\caption{Règle de décousure des \(2\)-liaisons d'un volume}
	\label{fig:formalisation-evenement-regles-jerboa-scission-rule-unsew-volume-a2}
\end{figure}

\begin{figure}[htb!]
	\centering
	\begin{subfigure}{.45\textwidth}
		\centering
		\includegraphics[height=100pt]{g-carte-unsew-volume-a2-G}
		\caption{G}
		\label{fig:formalisation-evenement-regles-jerboa-scission-unsew-volume-a2-G}
	\end{subfigure}
	\begin{subfigure}{.45\textwidth}
		\centering
		\includegraphics[height=100pt]{g-carte-unsew-volume-a2-H}
		\caption{H}
		\label{fig:formalisation-evenement-regles-jerboa-scission-unsew-volume-a2-H}
	\end{subfigure}
	\caption{Décousure des faces d'un volume}
	\label{fig:formalisation-evenement-regles-jerboa-scission-unsew-volume-a2}
\end{figure}

Étudions ce même cas, mais avec la règle de décousure des \(2\)-liaisons dans un
volume, c'est-à-dire une \(0,1,2\)-orbite. %
Dans cette règle, la scission est due à la suppression des \(2\)-arcs
implicites %
Ici, l'orbite est complète et son application déconnecte toutes les faces d'un
volume (Figure~\ref{fig:formalisation-evenement-regles-jerboa-scission-unsew-volume-a2}). %
Toutefois, si le volume correspond à une unique face, la scission sera tout de
même détectée. %
Dans une telle configuration, vérifier localement qu'il existe au moins une
\(2\)-liaisons qui n'est pas une boucle permet de confirmer qu'une scission a
bien eu lieu dans la transformation directe. %


\begin{figure}[htb!]
	\begin{subfigure}{.6\textwidth}
		\centering
		\includegraphics[width=\textwidth]{rule-cas-particulier-scission-1}
		\caption{Règle de décousure des sommets d'une arête}
		\label{fig:formalisation-evenement-regles-jerboa-scission-regle-cas-particulier-1}
	\end{subfigure}
	\begin{subfigure}{.4\textwidth}
		\centering
		\includegraphics[width=.8\textwidth]{rule-cas-particulier-scission-2}
		\caption{Règle de double décousure}
		\label{fig:formalisation-evenement-regles-jerboa-scission-regle-cas-particulier-2}
	\end{subfigure}
	\label{fig:formalisation-evenement-regles-jerboa-scission-regle-cas-particuliers}
\end{figure}

%% 4/ Deux composantes, une scission => scission entre les composantes
Étudions enfin le quatrième cas où une règle filtre deux composantes, mais où la
scission ne concerne que l'une des deux orbites dans la règle. %
Prenons une première règle qui découd les sommets d'une arête
(Figure~\ref{fig:formalisation-evenement-regles-jerboa-scission-regle-cas-particulier-1}). %
Les \(\orbit{0,1,3}\)-orbites incidentes à \(n0\) et \(n1\) sont toutes deux
incomplètes et l'orbite \(R\orbit{0,1,3}(n1)\) est scindée par son \(1\)-arc
explicite. %
Puisqu'aucune des orbites n'est complète, il est possible que les orbites
filtrées appartiennent à la même face. %
Dans ce cas, il y peut y avoir une scission entre les deux orbites filtrées
comme dans les
figures~\ref{fig:formalisation-evenement-regles-transfo-graphe-scission-cas-particulier-1-G}
et
\ref{fig:formalisation-evenement-regles-transfo-graphe-scission-cas-particulier-1-H}
(page~\pageref{fig:formalisation-evenement-regles-transfo-graphe-scission-cas-particulier-1-G}). %
% Inversement, si les deux orbites filtrées appartiennent à deux faces distinctes,
% alors il peut y avoir ou non une scission qui ne sépare pas les deux orbites
% filtrées.

Il s'agit du même principe pour la règle qui fait une double décousure
(Figure~\ref{fig:formalisation-evenement-regles-jerboa-scission-regle-cas-particulier-2}). %
Les \(\orbit{0,1,2}\)-orbites incidentes à \(n0\) et \(n1\) sont toutes deux incomplètes. %
Dans ce cas, il n'y a pas forcément de scission. %
Appliquée sur un ruban
(Figure~\ref{fig:formalisation-evenement-regles-transfo-graphe-scission-cas-particulier-2}),
une seule des deux scissions est nécessaire pour que le ruban soit effectivement
scindé comme nous pouvons le voir dans les figures. %
Ou alors, si les deux orbites désignent les extrémités du ruban, il ne se passe
rien, car les \(2\)-liaisons sont déjà des boucles. %

\begin{figure}[htb!]
	\centering
	\includegraphics[height=55pt]{rule-triangulate-face-compute-face-origin}
	\caption{Calcul de l'origine d'une face triangulée}
	\label{fig:formalisation-evenement-regles-jerboa-rule-triangulation-compute-face-origin}
\end{figure}

\begin{figure}[htb!]
	\begin{subfigure}{.45\textwidth}
		\centering
		\includegraphics[height=100pt]{g-carte-triangulate-face-compute-face-origin-G}
		\caption{\(G\)}
		\label{fig:formalisation-evenement-regles-jerboa-scission-triangulation-compute-face-origin-G}
	\end{subfigure}
	\begin{subfigure}{.45\textwidth}
		\centering
		\includegraphics[height=100pt]{g-carte-triangulate-face-compute-face-origin-H}
		\caption{\(H\)}
		\label{fig:formalisation-evenement-regles-jerboa-scission-triangulation-compute-face-origin-H}
	\end{subfigure}
	\caption{Origines des faces issues d'une triangulation}
	\label{fig:formalisation-evenement-regles-jerboa-scission-triangulation-compute-face-origin}
\end{figure}

La scission est le deuxième évènement pour lequel il est possible de calculer un
type d'origine. %
Calculons le type d'origine \(\orbit{\omega'}\) d'une face dans la règle de
triangulation
(Figure~\ref{fig:formalisation-evenement-regles-jerboa-rule-triangulation-compute-face-origin}). %
Dans cet exemple, le n\oe ud d'accroche contient trois dimensions,
\(\orbit{\omega'}\) contient donc trois dimensions. %
Dans l'orbite \(R\orbit{0,1,3}(n0)\) les réécritures sont les suivantes~:
\begin{itemize}
	\item le premier arc est soit réécrit en \(0\) soit supprimé, donc nous
	      ajoutons \(0\) à \(\orbit{\omega'}\). %
	\item le deuxième arc implicite est soit suppprimé soit réécrit en \(2\), donc
	      nous ajoutons \(2\) à \(\orbit{\omega'}\). %
	\item le troisième arc implicite est systématiquement réécrit en \(3\), donc nous ajoutons \(3\) à \(\orbit{\omega'}\).
\end{itemize}
Enfin, nous obtenons \(\orbit{\omega'} = \orbit{0,2,3}\). %
Ainsi chaque face issue de la triangulation, a pour origine une arête préservée
de la règle
(Figures~\ref{fig:formalisation-evenement-regles-jerboa-scission-triangulation-compute-face-origin-G}
et
\ref{fig:formalisation-evenement-regles-jerboa-scission-triangulation-compute-face-origin-H}). %

\subsection{Suppression}

\begin{figure}[htb!]
	\centering
	\includegraphics[height=50pt]{rule-dissolve-vertex-in-face-3D-no-color}
	\caption{Règle de suppression de l'étoile (d'un sommet et de ses arêtes incidentes)}
	\label{fig:formalisation-evenement-regles-jerboa-suppression-regle-dissolution-sommet}
\end{figure}


% La figure \textbf{fig:inst-rule-isolated-face-delete} représente une règle de
% suppression de face isolée où le nœud $n0$ dans $L$ filtre l'orbite de l'arête
% supprimée.

% \todo{Reporter corrections de contraction/dissolution d'un sommet en suppression
% 	de l'étoile}

Reprenons la règle supprime un sommet et son étoile
(Figure~\ref{fig:formalisation-evenement-regles-jerboa-suppression-regle-dissolution-sommet}). %
Dans la section~\ref{sec:formalisation-evenement-regles-transfo-graphe-suppression} (page~\pageref{sec:formalisation-evenement-regles-transfo-graphe-suppression}), nous
avons vu qu'un sommet est supprimé dans la transformation si et seulement si il
est supprimé dans la règle de transformation. %
Nous allons montrer que la suppression du sommet $L\orbit{1,2,3}(n2)$ dans la
règle Jerboa exprime la suppression d'un sommet dans la transformation
directe. %

\begin{proposition}[Suppression d'orbite]
	Soit une règle $r~: L \longrightarrow R$, un nœud d'accroche $h$ de $L$, un morphisme de
	filtrage $\underline{m}~: \underline{L} \rightarrow G$ sur une G-carte $G$, un type
	d'orbite $\orbit{\omega}$, $v$ un nœud de $L$ et $O = G\orbit{l(h)}(\underline{m}(h))$.

	L'orbite $\underline{L}\orbit{\omega}(v)$ est supprimée
	dans la règle Jerboa si et seulement si, pour
	tout brin $v'$ de $O$, l'orbite $G\orbit{\omega}(m((v,v')))$
	est supprimée dans la transformation directe $G \xlongrightarrow{r,\underline{m}} H$.
\end{proposition}

\begin{proof}
	Cette preuve est le symétrique de la preuve de la création (voir la proposition~\ref{prop:formalisation-evenement-regles-jerboa-creation}).
\end{proof}

% \begin{proof}
% 	Soit $v$ un nœud de $L$ tel que l'orbite $\underline{L}\orbit{\omega}(v)$ est
% 	supprimée par la règle, c'est-à-dire tel que
% 	$\underline{L}\orbit{\omega}(v) \cap \underline{R} = \emptyset$.

% 	Soit $w$ un nœud quelconque de l'orbite $\underline{L}\orbit{\omega}(v)$. $w$ est
% 	un nœud supprimé de la règle, c'est-à-dire de $\underline{L} \setminus \underline{R}$. Donc par construction de la règle
% 	instanciée, pour tout nœud $v'$ de l'orbite
% 	d'instanciation $O$, le n\oe ud $(w,v')$ est supprimé dans la règle instanciée
% 	$r(\underline{m})$.

% 	Donc toute orbite $L(O)\orbit{\omega}((v,v'))$ est supprimée dans la règle
% 	instanciée $r(\underline{m})$. %
%     Et donc, d'après le théorème de cohérence topologique~\ref{thm:preservation-coherence-topo-gcarte}, toutes les orbites
% 	$G\orbit{\omega}(m((v,v')))$ sont supprimées dans la transformation directe
% 	$G \xlongrightarrow{r,\underline{m}} H$, pour tous les nœuds $v'$ de l'orbite
% 	d'instanciation $O$.
% 	\\

% 	Réciproquement, supposons que toutes les orbites $G\orbit{\omega}(m((v,v')))$ sont
% 	supprimées dans la transformation directe $G \xlongrightarrow{r,\underline{m}} H$ pour tous
% 	les nœuds $v'$ de l'orbite d'instanciation $O$. Donc, d'après la proposition
% 	\ref{prop:base-rule-delete}, toutes les orbites $L(O)\orbit{\omega}((v,v'))$ pour
% 	tous les nœuds $v'$ de l'orbite d'instanciation $O$ sont des orbites
% 	supprimées dans la règle instanciée $r(\underline{m})$. Donc par construction
% 	de l'instanciation, $\underline{L}\orbit{\omega}(v)$ est une orbite supprimée dans
% 	la règle $r$.
% \end{proof}

\begin{figure}[htb!]
	\begin{subfigure}{.45\textwidth}
		\centering
		\includegraphics[height=100pt]{g-carte-triangulation-creation-H}
		\caption{\(G\)}
		\label{fig:formalisation-evenement-regles-jerboa-suppression-dissolution-sommet-G}
	\end{subfigure}
	\begin{subfigure}{.45\textwidth}
		\centering
		\includegraphics[height=100pt]{g-carte-triangulation-creation-G}
		\caption{\(H\)}
		\label{fig:formalisation-evenement-regles-jerboa-suppression-dissolution-sommet-H}
	\end{subfigure}
	\caption{Suppression de l'étoile sur un cube}
	\label{fig:formalisation-evenement-regles-jerboa-suppression-dissolution-sommet}
\end{figure}

D'après la condition d'arcs incidents du
théorème~\ref{thm:preservation-coherence-topo-regles}
(page~\pageref{thm:preservation-coherence-topo-regles}), une orbite supprimée
est nécessairement complète. %
Donc l'orbite à supprimer dans la règle instanciée est, elle aussi, complète. %
Donc l'orbite dans la G-carte est supprimée dans la transformation directe. %
Dans notre exemple
(Figure~\ref{fig:formalisation-evenement-regles-jerboa-suppression-regle-dissolution-sommet}),
l'orbite \(L\orbit{1,2,3}(n0)\) est complète donc sa suppression dans la règle
implique la suppression du sommet filtrée dans la transformation directe
(Figures~\ref{fig:formalisation-evenement-regles-jerboa-suppression-dissolution-sommet-G}\ref{fig:formalisation-evenement-regles-jerboa-suppression-dissolution-sommet-H}). %

\subsection{Fusion}

Reprenons la règle qui supprime une étoile
(Figure~\ref{fig:formalisation-evenement-regles-jerboa-suppression-regle-dissolution-sommet}). %
Dans la section~\ref{sec:formalisation-evenement-regles-transfo-graphe-fusion}
(page~\pageref{sec:formalisation-evenement-regles-transfo-graphe-fusion}), nous
avons vu que les face incidente au sommet supprimé sont fusionnées dans la
transformation directe si elles sont transformées dans la règle instanciée. %
Nous allons montrer que la fusion des \(\orbit{0,1,3}\)-orbites dans la règle
implique une fusion de faces dans la transformation directe. %

\begin{proposition}[Fusion d'orbites]
	Soit une règle Jerboa $r~: L \longrightarrow R$, un nœud d'accroche $h$ de $L$, un morphisme
	de filtrage $\underline{m}~: \underline{L} \rightarrow G$ sur une G-carte $G$, un type
	d'orbite $\orbit{\omega}$, deux n\oe uds préservés $v_1$ et $v_2$ de
	$\underline{K} = \underline{L} \cap \underline{R}$ et l'orbite d'instanciation
	$O = G\orbit{l(h)}(\underline{m}(h))$. %

	$\underline{L}\orbit{\omega}(v_1)$ est fusionnée dans la règle~:

	\begin{itemize}
		\item si la fusion est due à des arcs explicites, c'est-à-dire si
		      $\underline{L}\orbit{\omega}(v_1)$ est fusionnée avec
		      $\underline{L}\orbit{\omega}(v_2)$ dans $\underline{R}$~;

		\item si la fusion est due à des arcs implicites, c'est-à-dire si pour tout
		      nœud $v$ de $\underline{L}\orbit{\omega}(v_1)$, l'étiquette $l(v)$ a pour
		      $k$-ème position un "\_" ou une dimension qui n'appartient pas à
		      $\orbit{\omega}$, et qu'il existe un nœud $w$ de
		      $\underline{R}\orbit{\omega}(v_1)$ dont l'étiquette $l(w)$ a pour $k$-ème
		      position une dimension de $\orbit{\omega}$.
	\end{itemize}

	\begin{enumerate}
		\item Si $\underline{L}\orbit{\omega}(v_1)$ et $\underline{L}\orbit{\omega}(v_2)$ sont
		      fusionnées en une orbite $\underline{R}\orbit{\omega}(v_1)$ dans la règle
		      du à des arcs explicites. Et si l'orbite $\underline{L}\orbit{\omega}(v_1)$
		      est complète. %
		      Alors, pour tout brin $v'$ de $O$, les orbites
		      $G\orbit{\omega}(m((v_1,v')))$ et $G\orbit{\omega}(m((v_2,v')))$ sont fusionnées
		      en une orbite $H\orbit{\omega}(m'((v_1,v')))$ dans la transformation
		      directe $G \xlongrightarrow{r,\underline{m}} H$.

		\item Si les orbites $\underline{L}\orbit{\omega}(v_1)$ et
		      $\underline{L}\orbit{\omega}(v_2)$ sont fusionnées en une orbite
		      $\underline{R}\orbit{\omega}(v_1)$ dans la règle du à des arcs explicites.
		      Et si les orbites $\underline{L}\orbit{\omega}(v_1)$ et
		      $\underline{L}\orbit{\omega}(v_2)$ sont toutes deux incomplètes. %
		      Alors, pour chaque brin $v'$ de $O$, soit $G\orbit{\omega}(m((v_1,v')))$ et
		      $G\orbit{\omega}(m((v_2,v')))$ sont fusionnées en une orbite
		      $H\orbit{\omega}(m'((v_1,v')))$, soit elles sont la même orbite modifiée
		      dans la transformation directe
		      $G \xlongrightarrow{r, \underline{m}} H$.

		\item Si l'orbite $\underline{L}\orbit{\omega}(v_1)$ est fusionnées dans la règle
		      du au $k$-ème arc implicite. Alors pour tout brin $v_1'$ de $O$ et
		      $v_2'$ le $i$-voisin de $v_1'$ dans $O$ avec $i$ la dimension à la
		      $k$-ème position de $l(h)$, $G\orbit{\omega}(m((v_1,v_1')))$ et
		      $G\orbit{\omega}(m((v_1,v_2')))$ sont soit fusionnées en une orbite
		      $H\orbit{\omega}(m'((v_1,v_1')))$, soit c'est une même orbite modifiée
		      $H\orbit{\omega}(m'((v_1,v_1'))) = H\orbit{\omega}(m'((v_1,v_2')))$ dans la
		      transformation directe $G \xlongrightarrow{r,\underline{m}} H$
		\item Réciproquement,

		      si les orbites $L\orbit{\omega}(m((v_1, v'_1))$ et
		      $L\orbit{\omega}(m((v_2, v'_2))$ sont fusionnées en une orbite
		      $H\orbit{\omega}(m'((v_1, v'_1)))$, alors il existe deux n\oe uds préservés
		      $v_3$ et $v_4$ (de $\underline{K} = \underline{L} \cap \underline{R}$)
		      tel que $m((v_3, v'_1)) \in H\orbit{\omega}(m'((v_1, v'_1))$ et
		      $m((v_4, v'_1)) \in H\orbit{\omega}(m'((v_1, v'_1))$ et
		      $\underline{L}\orbit{\omega}(v_3)$ et $\underline{L}\orbit{\omega}(v_4)$ sont
		      fusionnées en
		      $\underline{R}\orbit{\omega}(v_3) = \underline{R}\orbit{\omega}(v_4) $. %
	\end{enumerate}
\end{proposition}

\begin{proof}
	Cette preuve est la symétrique de celle de la scission (voir la proposition~\ref{prop:formalisation-evenement-regles-jerboa-scission}). %
\end{proof}

% \begin{proof}
% 	Soit deux nœuds préservés $v_1$ et $v_2$ de $L$ et $R$
% 	\begin{enumerate}
% 		\item Supposons que $\underline{L}\orbit{\omega}(v_1)$ et
% 		      $\underline{L}\orbit{\omega}(v_2)$ sont fusionnées en une orbite
% 		      $\underline{R}\orbit{\omega}(v_1)$ dans la règle du à des arcs
% 		      explicites. Alors, par définition, il n'existe pas de chemin
% 		      $e_1 e_2 \dots e_j$ dans $L$ de source $v_1$, de cible $v_2$ et
% 		      étiqueté par $i_1 i_2 \dots i_j$ sur $\orbit{\omega}$, mais il existe un
% 		      chemin $f_1 f_2 \dots f_k$ dans $R$ de source $v_1$, de cible $v_2$
% 		      et étiqueté sur $\orbit{\omega}$.

% 		      Donc dans la règle instanciée, pour tout nœud $v'$ de l'orbite d'instanciation $O = G\orbit{l(h)}(\underline{m}(h))$, il n'existe pas de chemin $(e_1, v')(e_2, v')\dots(e_j,v')$ dans le
% 		      motif gauche de la règle instanciée $L(O)$ dont la source est $(s(e_1),v') = (v_1,v')$, la cible est $(t(e_j),v') = (v_2,v')$ et étiqueté par le même mot $i_1 i_2 \dots i_j$ sur $\orbit{\omega}$.

% 		      Supposons que pour tout nœud $v'$ de l'orbite d'instanciation $O$, il n'existe pas de chemin $(f_1,v')(f_2,v')\dots(f_k,v')$ dans le motif droit de la règle instanciée $R(O)$ dont la source
% 		      est $(s(f_1),v') = (v_1,v')$, la cible est $(t(f_k),v') = (v_2, v')$ et étiqueté sur $\orbit{\omega}$.
% 		      Alors, par construction de l'instanciation, le chemin $f_1 f_2 \dots f_k$ dans $R$ de source $v_1$,
% 		      de cible $v_2$ et étiqueté sur $\orbit{\omega}$ n'existe pas dans la règle de départ.
% 		      Ce qui est contradictoire avec l'hypothèse. Donc il existe un chemin $(f_1,v')(f_2,v')\dots(f_k,v')$ dans le
% 		      motif droit de la règle instanciée ayant les propriétés attendues.
% 		      Donc $R(O)\orbit{\omega}((v_1,v'))$ et $R(O)\orbit{\omega}((v_2,v'))$ sont bien une seule et même $\orbit{\omega}$-orbite.

% 		      Donc d'après la proposition \ref{prop:base-rule-merge}, si $\underline{L}\orbit{\omega}(v_1)$ et donc $L(O)\orbit{\omega}((v_1,v'))$ est entièrement filtrée, alors les orbites $G\orbit{\omega}(m((v_1,v')))$ et
% 		      $G\orbit{\omega}(m((v_2,v')))$ sont fusionnées en une orbite $H\orbit{\omega}(m'((v_1,v')))$ dans la transformation directe $G \xlongrightarrow{r,\underline{m}} H$ pour tout nœud $v'$ de l'orbite
% 		      d'instanciation $O$.

% 		\item Supposons que $\underline{L}\orbit{\omega}(v_1)$ et
% 		      $\underline{L}\orbit{\omega}(v_2)$ sont fusionnées en une orbite
% 		      $\underline{R}\orbit{\omega}(v_1)$ dans la règle du à des arcs
% 		      explicites.

% 		      Pour les mêmes raisons que dans le cas $1$, $L(O)\orbit{\omega}((v_1,v'))$ et $L(O)\orbit{\omega}((v_2,v'))$ sont deux $\orbit{\omega}$-orbites distinctes dans la règle instanciée, pour tout nœud $v'$ de l'orbite d'instanciation $O$.

% 		      Donc d'après la proposition \ref{prop:base-rule-merge}, si $\underline{L}\orbit{\omega}(v_1)$ et donc $L(O)\orbit{\omega}((v_1,v'))$ est partiellement filtrée, alors les orbites
% 		      $G\orbit{\omega}(m((v_1,v')))$ et $G\orbit{\omega}(m((v_2,v')))$ sont soit fusionnées en une orbite $H\orbit{\omega}(m'((v_1,v')))$, soit modifiées, dans la transformation directe
% 		      $G \xlongrightarrow{r,\underline{m}} H$ pour tout nœud $v'$ de l'orbite d'instanciation $O$.

% 		\item Supposons que $\underline{L}\orbit{\omega}(v_1)$ est fusionnée due au
% 		      $k$-ème arc implicite. Soit $v_1'$ un nœud de l'orbite
% 		      d'instanciation $O = G\orbit{l(h)}(\underline{m}(h))$ et
% 		      $v_2'$ le $i$-voisin de $v_1'$ dans $O$, avec $i$ la $k$-ème
% 		      dimension de l'étiquette $l(h)$ du nœud d'accroche $h$ de $L$.

% 		      Par hypothèse, tous les n\oe uds de $\underline{L}\orbit{\omega}(v_1)$ ont une $k$-ième dimension soit supprimée, soit renommée hors de $\orbit{\omega}$.
% 		      Au contraire, il existe un nœud $v$ de $\underline{R}\orbit{\omega}(v_1)$ dont la $k$-ième dimension est renommée en $j$ dans $\orbit{\omega}$. Donc les nœuds $(v, v'_1)$ et $(v, v'_2)$ sont $j$-voisins dans l'instanciation $R(O)$, donc $(v, v'_1)$ et $(v, v'_2)$ sont dans la même $\orbit{\omega}$-orbite dans $R(O)$, et donc $(v_1, v'_1)$ et $(v_2, v'_2)$ aussi.

% 		      Nous avons plusieurs cas :
% 		      \begin{itemize}
% 			      \item si $(v_1, v'_1)$ et $(v_2, v'_2)$ sont deux $\orbit{\omega}$-orbites différentes de $L(O)$, d'après la proposition \ref{prop:base-rule_split},
% 			            selon les cas, $G\orbit{\omega}(m((v_1, v'_1)))$ et $G\orbit{\omega}(m((v_1,v'_2)))$ sont dans deux orbites différentes ou non de $G$ ;
% 			            autrement dit $G\orbit{\omega}(m((v_1, v'_1))$ et $G\orbit{\omega}(m((v_2, v'_2))$ sont fusionnées en une orbite $H\orbit{\omega}(m'((v_1, v'_1)))$ ou l'unique orbite $G\orbit{\omega}(m((v_1, v'_1)) = G\orbit{\omega}(m((v_2, v'_2))$ est juste modifiée par la transformation directe $G \xlongrightarrow{r, \underline{m}} H$ ;

% 			      \item si au contraire $(v_1, v'_1)$ et $(v_2, v'_2)$ sont dans la même $\orbit{\omega}$-orbites de $L(O)$ (car il existe un chemin $e_1 e_2 ... e_l$ de $O$ de source $v'_1$ et de cible $v'_2$ qui ne comprend pas l'étiquette $j$, ce chemin est réécrit un chemin $(w, e_1) (w,e_2) ... (w, e_l)$ de $L(O)$ dont toutes les étiquettes sont dans $\orbit{\omega}$),
% 			            d'après la proposition \ref{prop:base-rule-modify}, alors $G\orbit{\omega}(m((v_1, v'_1))$ est modifiée par la transformation directe $G \xlongrightarrow{r, \underline{m}} H$.
% 		      \end{itemize}
% 	\end{enumerate}
% \end{proof}

La fusion est l'évènement symétrique de la scission. %
Ainsi la méthode de détection d'une fusion est la même. %

\begin{figure}[htb!]
	\centering
	\includegraphics[height=50pt]{rule-sewA3}
	\caption{Règle de coutures de faces de volumes}
	\label{fig:formalisation-evenement-regles-jerboa-fusion-regle-couture-faces-a3}
\end{figure}

\begin{figure}[htbp!]
	\centering
	\begin{subfigure}{.45\textwidth}
		\centering
		\includegraphics[width=.85\textwidth]{g-carte-decousure-3-H}
		\caption{$G$}
		\label{fig:formalisation-evenements-regles-jerboa-evenement-fusion-explicite-g}
	\end{subfigure}
	\begin{subfigure}{.45\textwidth}
		\centering
		\includegraphics[width=.85\textwidth]{g-carte-decousure-3-G}
		\caption{$H$}
		\label{fig:formalisation-evenement-regles-jerboa-evenement-fusion-explicite-face-h}
	\end{subfigure}
	\caption{Décousure des $3$-liaisons dans une face}
	\label{fig:formalisation-evenement-regle-evenement-fusion-explicite}
\end{figure}

Pour illustrer le premier cas, une fusion explicite où il existe au moins une
\(\orbit{\omega}\)-orbite complète, étudions la règle de couture de deux faces par
des \(3\)-arcs. %
Dans cette règle, les orbites faces sont toutes complètes, leur connexion par le
\(3\)-arc provoque leur fusion dans la règle. %
Donc, dans la transformation directe, deux faces, entièrement filtrées, sont
connectés par des \(3\)-liaisons et donc fusionnées en une
\(\orbit{0,1,3}\)-orbite
(Figures~\ref{fig:formalisation-evenements-regles-jerboa-evenement-fusion-explicite-g}
et
\ref{fig:formalisation-evenement-regles-jerboa-evenement-fusion-explicite-face-h}). %

\begin{figure}[htb!]
	\centering
	\includegraphics[height=50pt]{rule-delete-edge-no-color}
	\caption{Règle de suppression d'une arête}
	\label{fig:formalisation-evenement-regles-jerboa-fusion-regle-suppression-arete}
\end{figure}

Pour le deuxième cas, une fusion où il n'existe pas de \(\orbit{\omega}\)-complète,
prenons la règle fusionne deux face incidentes à une même arête. %
Dans cette règle, aucun n\oe ud de l'orbite \(L\orbit{0,1,3}(n0)\) n'a de
deuxième arc implicite étiqueté dans \(\orbit{0,1,3}\), mais il y en a un dans
\(R\orbit{0,1,3}(n1)\). %
Nous détectons donc une fusion par des \(1\)-arcs implicites dans la règle
néanmoins les orbites sont toutes deux incomplètes. %
Il est donc nécessaire de vérifier si la règle provoque une scission dans la
transformation directe. %



\subsection{Non-changement}

% Dans l'insertion de sommet figure \textbf{fig:rule-insert-vertex}, le nœud $n0$
% (en noir) filtre partiellement les orbites sommet aux extrémités de l'arête sur
% laquelle est inséré le nouveau sommet. Ces orbites sont inchangées par
% l'insertion du nouveau sommet.

La règle de triangulation
(Figure~\ref{fig:formalisation-evenement-regles-jerboa-rule-face-triangulation} page~\pageref{fig:formalisation-evenement-regles-jerboa-rule-face-triangulation}),
réécrit les étiquettes du n\oe ud $n0$, qui étaient $\orbit{0,\_,3}$ dans $L$ et
deviennent $\orbit{0,1,3}$ dans $R$, l'étiquette $0$ est préservée. %
Le n\oe ud $n0$ filtre ainsi des arêtes qui restent inchangées suite à la
réécriture des étiquettes du n\oe ud. %

\begin{proposition}[Non-changement d'orbite]
	Soit une règle Jerboa $r~: L \longrightarrow R$, un nœud d'accroche $h$ de $L$,
	un morphisme de filtrage $\underline{m}~: \underline{L} \rightarrow G$ sur une G-carte
	$G$, un type d'orbite $\orbit{\omega}$, $v$ un nœud de $L$ et l'orbite d'instanciation $O = G\orbit{l(h)}(\underline{m}(h))$.

	L'orbite $\underline{L}\orbit{\omega}(v)$ est inchangée dans
	la règle si et seulement si, pour tout brin $v'$ de $O$, l'orbite $G\orbit{\omega}(m((v,v')))$ est
	inchangée dans la transformation directe $G \xlongrightarrow{r,\underline{m}} H$.
\end{proposition}

\begin{proof}
	Supposons que $\underline{L}\orbit{\omega}(v)$ est inchangée dans la règle Jerboa.

	Soit $w$ un nœud quelconque de l'orbite $\underline{L}\orbit{\omega}(v)$. %
	Donc $w$ est un	nœud inchangé de la règle de $\underline{K} = \underline{L} \cap \underline{R}$. Et donc, par construction de la règle
	instanciée, pour tout brin $v'$ de $O$, le n\oe ud $(w,v')$ est inchangé dans la règle instanciée	$r(\underline{m})$. %
	Or l'orbite $G\orbit{\omega}(m((v,v')))$ n'est pas
	nécessairement complète. %
	Par définition, la partie non filtrée est aussi inchangée. %
	Donc $L(O)\orbit{\omega}((v,v'))$ est inchangée dans la règle instanciée. %
	Et donc, d'après la proposition
	\ref{prop:base-rule-unchange}, pour tout brin $v'$ de $O$, l'orbite $G\orbit{\omega}(m((v,v')))$ est
	inchangée dans la transformation directe $G \xlongrightarrow{r,\underline{m}} H$.

	Réciproquement, supposons que pour tout brin $v'$ de $O$ l'orbite $G\orbit{\omega}(m((v,v')))$ est
	inchangée dans la transformation directe $G \xlongrightarrow{r, \underline{m}} H$. %
	Donc d'après la proposition	\ref{prop:base-rule-unchange}, l'orbite $L(O)\orbit{\omega}((v,v'))$  est inchangée
	dans la règle instanciée. %
	Donc par construction de l'instanciation, l'orbite $\underline{L}\orbit{\omega}(v)$ est inchangée dans la règle Jerboa. %
\end{proof}

\maxime{Instanciation + description}

\begin{remarque}
	L'orbite \(G\orbit{\omega}(b)\), où \(b\) est un brin non filtré de \(G\), n'est \emph{pas affectée} par l'application de $r$ si
	$G\orbit{\omega}(b) \cap \underline{m}(L) = \emptyset$ et \emph{inchangée}
	sinon.
\end{remarque}

\todo{illustration + commentaire}

% \subsection{Sans effet \todo{redéfinir comme une remarque dans non-changement}}

% \maxime{appliquer la règle d'insertion de sommet sur un carré}

% Dans la figure ? le brin $a$ de $G$ n'est pas filtré, mais l'orbite sommet $G\orbit{1,2,3}(a)$ est partiellement filtrée par la règle \maxime{à continuer en fonction du dessin} […] donc l'orbite sommet $G\orbit{1,2,3}(a)$ est inchangée.
% Le brin $b$ de $G$ n'est pas filtré et l'orbite $G\orbit{1,2,3}(b)$ n'est pas filtrée non plus, donc l'orbite sommet $G\orbit{1,2,3}(b)$ n'est pas affectée par l'application de la règle.

% \todo{référence en avant : on sera amené à faire la distinction entre les orbites au moins partiellement filtrées et ce qui n'est pas filtré du tout.}

% \begin{definition}[Orbite non affectée]
% 	Soit une règle $r~: L \longrightarrow R$, un nœud d'accroche $h$ de $L$, un morphisme de filtrage
% 	$\underline{m}~: \underline{L} \rightarrow G$ sur une G-carte $G$, et un type d'orbite
% 	$\orbit{\omega}$.

% 	Soit $b$ un brin de $G$ non filtré. L'orbite $G\orbit{\omega}(b)$ n'est \emph{pas affectée} par l'application de $r$ si $G\orbit{\omega}(b) \cap \underline{m}(L) = \emptyset$ et \emph{inchangée} sinon.

% \end{definition}

\subsection{Modification}

\todo{la modification ce sont tous les autres cas (augmentation, réduction, réécriture des liens topologiques [par exemple transformation d'une face en un sommet,…],…}

\begin{proposition}[Modification d'orbite]
	Soit une règle Jerboa $r~: L \longrightarrow R$, un nœud d'accroche $h$ de $L$, un morphisme de filtrage $\underline{m}~: \underline{L} \rightarrow G$ sur une G-carte $G$, un type d'orbite $\orbit{\omega}$, $v$ un nœud préservé de $\underline{K} = \underline{L} \cap \underline{R}$ et l'orbite d'instanciation	$O = G\orbit{l(h)}(\underline{m}(h))$.

	L'orbite $\underline{L}\orbit{\omega}(v)$ est modifiée dans la règle si et seulement si~:
	\begin{itemize}
		\item si la modification est due à des arcs explicites, alors~:
		      \begin{itemize}
			      \item $\underline{L}\orbit{\omega}(v) \neq \underline{R}\orbit{\omega}(v)$,
			      \item pour tout n\oe ud préservé $w$ de $\underline{K} = \underline{L} \cap \underline{R}$, $w \in \underline{L}\orbit{\omega}(v)$ si et seulement si $w \in \underline{R}\orbit{\omega}(v)$,
			      \item pour toute position d'étiquette $k$, il existe un n\oe ud $w \in \underline{L}\orbit{\omega}(v)$ dont la $k$-ème étiquette de $l(w)$ est une dimension de $\orbit{\omega}$ si et seulement si il existe un n\oe ud $w' \in \underline{R}\orbit{\omega}(v)$ dont la $k$-ème étiquette de $l(w')$ est une dimension de $\orbit{\omega}$~;
		      \end{itemize}
		\item si la modification est due à des arcs implicites, alors~:
		      \begin{itemize}
			      \item $\underline{L}\orbit{\omega}(v) \neq \underline{R}\orbit{\omega}(v)$,
			      \item il existe un n\oe ud préservé $w \in \underline{L}\orbit{\omega}(v)$ tel que l'étiquette de $w$ est différente dans $L$ et $R$,
			      \item pour toute position d'étiquette $k$, il existe un n\oe ud $w \in \underline{L}\orbit{\omega}(v)$ dont la $k$-ème étiquette de $l(w)$ est une dimension de $\orbit{\omega}$ si et seulement si il existe un n\oe ud $w' \in \underline{R}\orbit{\omega}(v)$ dont la $k$-ème étiquette de $l(w')$ est une dimension de $\orbit{\omega}$~;
		      \end{itemize}
	\end{itemize}


	Si l'orbite $\underline{L}\orbit{\omega}(v)$ est modifiée dans la règle Jerboa, alors, pour tout brin $v'$ de $O$,
	$G\orbit{\omega}(m((v,v')))$ est modifiée dans la transformation directe
	$G \xlongrightarrow{r,\underline{m}} H$. %
	Réciproquement, si 	$G\orbit{\omega}(m((v,v')))$ est modifiée dans la transformation directe
	$G \xlongrightarrow{r,\underline{m}} H$, alors $\underline{L}\orbit{\omega}(v)$ est modifiée, scindée ou fusionnée dans la règle. %
\end{proposition}

\begin{proof}
	Supposons que l'orbite $\underline{L}\orbit{\omega}(v)$ est modifiée dans la règle due à des arcs explicites. %
	Montrons que pour tout brin $v'$ de l'orbite d'instanciation $O$, l'orbite $L(O)\orbit{\omega}((v, v'))$ est modifié dans la règle instanciée. %

	D'après le premier point, les orbites $L(O)\orbit{\omega}((v, v'))$ et $R(O)\orbit{\omega}((v, v'))$ sont différentes. %
	Elles ne peuvent donc pas être inchangées. %
	Il nous reste donc à montrer qu'elles ne sont ni scindées ni fusionnées. %

	Soit un nœud préservé $w$ de $\underline{K} = \underline{L} \cap \underline{R}$ et deux brins $v'_1$ et $v'_2$ de l'orbite d'instanciation $O$. %
	Supposons que $(w, v'_1)$ et $(v, v'_2)$ soient dans la même $\orbit{\omega}$-orbite. %
	Alors, il existe un chemin $c(O)$ dans $L(O)$ étiqueté sur $\orbit{\omega}$ de source $(w, v'_1)$ et de cible $(v, v'_2)$. %
	Les arcs de ce chemin issus d'arcs explicites de $L$, forment un chemin $c$ de $L$ entre $w$ et $v$ qui s'instancie en un chemin entre $(w, v'_1)$ et $(v, v'_1)$ dans $L(O)$. %
	Or par hypothèse, il existe un chemin issu d'arcs explicites de $R$, qui forment un chemin de $R$ entre $w$ et $v$ qui s'instancie en un chemin entre $(w, v'_1)$ et $(v, v'_1)$ dans $R(O)$. %
	Les arcs de $c(O)$ issus d'arcs implicites de $L$ forment un chemin de $O$ entre $v'_1$ et $v'_2$. %
	Or pour chacun de ces arcs de $v'_i$ à $v'_j$ correspond un arc implicite de $L\orbit{\omega}(v)$, et il existe un arc implicite dans $R\orbit{\omega}(v)$ étiqueté sur $\orbit{\omega}$, éventuellement sur un autre nœud $w' \in \underline{R}\orbit{\omega}(v)$. %
	On peut donc construire un chemin dans $R(O)$ qui utilise l'instanciation d'arcs explicites de $(w, v'_i)$ à $(w', v'_i)$, puis l'instanciation de l'arc implicite pour passer à $(w', v'_j)$, puis à nouveau l'instanciation d'arcs explicites pour passer à $(w, v'_j)$. %
	On peut donc construire un chemin de $(w, v'_1)$ à $(w, v'_2)$ dans $R(O)$. %
	Nous avons donc construit un chemin de $(w, v'_1)$ à $(v, v'_2)$ dans $R(O)$ étiqueté sur $\orbit{\omega}$.

	La construction inverse d'un chemin de $R(O)$ vers un chemin de $L(O)$ peut se faire de la même manière. %
	Donc pour tout nœud $v'$ de l'orbite d'instanciation $O$, l'orbite $L(O)\orbit{\omega}((v, v'))$ est modifiée dans la règle instanciée.
	\\ %

	Supposons que l'orbite $\underline{L}\orbit{\omega}(v)$ est modifiée dans la règle due à des arcs implicites. %
	Montrons que pour tout nœud $v'$ de l'orbite d'instanciation $O$, l'orbite $L(O)\orbit{\omega}((v, v'))$ est modifiée dans la règle instanciée. %

	Certains arcs implicites des orbites considérées sont différentes dans $L$ et $R$, donc pour tout nœud $v'$ de l'orbite d'instanciation $O$, les orbites $L(O)\orbit{\omega}((v, v'))$ et $R(O)\orbit{\omega}((v, v'))$ sont différentes. %
	Il nous reste donc à montrer qu'il n'y a ni scission ni fusion. %

	Cela se fait de la même manière que dans le cas explicite, avec la partie de chemin issue des arcs explicites identiques dans $L(O)$ et $R(O)$, puisque les structures $\underline{L}\orbit{\omega}(v)$ et $\underline{R}\orbit{\omega}(v)$ sont identiques. %
	\\ %

	Donc, d'après la proposition \ref{prop:base-rule-modify}, pour tout brin $v'$ de $O$,
	$G\orbit{\omega}(m((v,v')))$ est modifiée dans la transformation directe
	$G \xlongrightarrow{r,\underline{m}} H$. %
	\\ %

	Réciproquement, soit un brin $v'$ de $O$, supposons que $G\orbit{\omega}(m((v,v')))$ est modifiée dans la transformation directe $G \xlongrightarrow{r,\underline{m}} H$. %
	D'après les propositions précédentes, l'orbite $L\orbit{\omega}(v)$ ne peut être ni inchangée, ni scindée, ni fusionnée (et bien sûr ni crée, ni supprimée). %
	Elle est donc nécessairement modifiée.
\end{proof}


\subsection{Exemples de détections d'évènements}

\begin{table}[htbp!]
	\caption{\label{tab:formalisation-evenement-regle-jerboa-tableau-detection}
		Extrait des évènements détectés lors de l'application de la règle de triangulation
		(Figure~\ref{fig:formalisation-evenement-regles-jerboa-rule-face-triangulation})
		sur la face d'un cube (Figure~\ref{fig:formalisation-evenement-regles-jerboa-inst-quad-triangulation})}

	\begin{tabular}{|p{0.1\textwidth}|p{0.846\textwidth}|}
		\hline
		\centering Type d'orbite   & Évènements détectés                                                                                                                       \\
		\hline
		\centering $\orbit{0}$     & Création $H\orbit{0}(a_{1})$
		\par Création~$H\orbit{0}(b_{1})$ \ldots
		\vspace{0.1 cm}
		\par Non-changement~$H\orbit{0}(a_{0})$
		\par Non-changement~$H\orbit{0}(c_{0})$~\ldots                                                                                                                         \\
		\hline
		\centering $\orbit{1}$     & Création $H\orbit{0}(a_{2})$
		\par Création~$H\orbit{1}(c_{2})$ \ldots
		\vspace{0.1 cm}
		\par Scission~$G\orbit{1}(a_{0})$ $\rightarrow\{H\orbit{1}(a_{0})$, $H\orbit{1}(h_{0})\}$
		\par Scission~$G\orbit{1}(c_{0})$ $\rightarrow \{H\orbit{1}(c_{0})$, $H\orbit{1}(b_{0})\}$~\ldots                                                                      \\
		\hline
		\centering $\orbit{2}$     & Création $H\orbit{2}(a_{1})$
		\par Création~$H\orbit{2}(a_{2})$
		\par Création~$H\orbit{2}(b_{1})$
		\par Création~$H\orbit{2}(b_{2})$ \ldots
		\vspace{0.1 cm}
		\par Non-changement~$H\orbit{2}(a_{0})$
		\par Non-changement~$H\orbit{2}(b_{0})$
		\par Non-changement~$H\orbit{2}(c_{0})$
		\par Non-changement~$H\orbit{2}(d_{0})$~\ldots                                                                                                                         \\

		\hline
		\centering $\orbit{1,2}$   & Modification $H\orbit{1,2}(a_{0})$
		\par Modification~$H\orbit{1,2}(b_{0})$\ldots
		\vspace{0.1 cm}
		\par Création~$H\orbit{1,2}(a_{2})$                                                                                                                                    \\
		\hline
		\centering $\orbit{0,1}$   & Scission~$G\orbit{0,1}(a_{0})$ $\rightarrow \{H\orbit{0,1}(a_{0})$, $H\orbit{0,1}(c_{0})$, $H\orbit{0,1}(e_{0})$, $H\orbit{0,1}(g_{0})\}$ \\
		\hline
		\centering $\orbit{0,2}$   & Création $H\orbit{0,2}(a_{1})$
		\par Création~$H\orbit{0,2}(b_{1})$\ldots
		\vspace{0.1 cm}
		\par Non-changement~$H\orbit{0,2}(a_{0})$
		\par Non-changement~$H\orbit{0,2}(c_{0})$~\ldots                                                                                                                       \\
		\hline
		\centering $\orbit{0,1,2}$ & Modification~$H\orbit{0,1,2}(a_{0})$                                                                                                      \\
		\hline
	\end{tabular}
	\normalsize

\end{table}

%%% Local Variables:
%%% mode: latex
%%% TeX-master: "../../../main"
%%% End:
